\naslov{Sistemi linearnih NDE}

Naj bodo podane funkcije $a_{ij} : [a,b] \to \R$ in $b_k: [a,b] \to \R$, ki so
na $[a,b]$ omejene.
Sistem NDE prvega reda s koeficienti $a_{ij}(t)$ in desno stranjo $b_k(t)$ je
sistem
\begin{align*}
  \dot{x}_1 &= a_{11} x_1 + \ldots + a_{1n} x_n \\
   & \vdots \\
  \dot{x}_n &= a_{n1} x_1 + \ldots + a_{nn} x_n \\
\end{align*}
Naj bo matrika $A(t)$ podana s koeficienti $a_{ij}$ in $b$ vektor podan s
komponentami $b_k$.
Za $x = [x_1 \ldots x_n]^T$ sistem zapišemo kot $\dot{x} = Ax + b$.
Če je $b = 0$ pravimo, da je sistem \pojem{homogen}.

\begin{izrek}
  Če je $A:[a,b] \to \R^{n \times n}$ zvezna in omejena, je množica rešitev
  homogenega sistema $\dot{x} = Ax$ je $n$-dimenzionalen vektorski prostor v
  prostoru $\zvodv{1}{[a,b]}$.
\end{izrek}

\begin{proof}
  Naj bo $R$ prostor rešitev.
  Najprej moramo pokazati, da je $R$ vektorski podprostor.
  Za vsak par rešitev $x_1, x_2$ in $\alpha, \beta \in \R$ velja
  \[
	A(\alpha x_1 + \beta x_2) = \alpha A x_1 + \beta A x_2 = \partial_t (\alpha
	x_1 + \beta x_2).
  \]
  Naj bo $t_0 \in [a,b]$.
  Po eksistenčnem izreku obstaja rešitev vsakega začetnega problema $\dot{x} =
  Ax, x(t_0) = c \in \R^n$.
  Naj bo $e_1, \ldots, e_n$ kanonična baza v $\R^n$.
  Obstajajo torej rešitve $x_i$ začetnih problemov $\dot{x} = Ax, x(t_0) = e_i$.
  Dokazati moramo še, da so $x_i$ tudi globalne rešitve; ta del pustimo za
  kasneje, preostanek dokaza je lokalen.

  Trdimo, da je $\{x_i\}_i$ baza $R$.
  Linearna neodvisnost je trivialna.
  Naj bo $x$ rešitev sistema in $c = x(t_0)$.
  Vektor $c$ razvijemo po bazi $e_i$ v $c = \alpha_1 e_1 + \ldots + \alpha_n
  e_n$ in definiramo
  \[
	\tilde{x} = \alpha_1 x_1 + \ldots + \alpha_n x_n.
  \]
  Ker sta tako $x$ kot $\tilde{x}$ rešitvi začetnega problema $\dot{x} = Ax,
  x(t_0) = c$, sta po eksistenčnem izreku enaki.
\end{proof}

\vprasanje{Kaj je množica rešitev homogenega sistema linearnih NDE? Dokaži.}

\begin{definicija}
  \pojem{Fundamentalna matrika} sistema $\dot{x} = Ax$ je matrika
  \[
	\phi(t,t_0)
	=
	\begin{bmatrix}
	  x_{11}(t) & \ldots & x_{1n}(t) \\
	  \vdots & \ddots & \vdots \\
	  x_{n1}(t) & \ldots & x_{nn}(t)
	\end{bmatrix},
  \]
  v kateri je $i$-ti stolpec enak $i$-ti rešitvi iz dokaza.
\end{definicija}

Za matriko $\phi(t,t_0)$ velja $\phi(t_0,t_0) = I$.
Naj bo $x$ rešitev začetnega problema $\dot{x} = Ax, x(t_0) = c$.
Potem velja $x = \phi(t,t_0)c$.

\begin{trditev}
  Za vsak $t$ veja $\det \phi(t,t_0) \ne 0$.
\end{trditev}

\begin{proof}
  Recimo, da obstaja $t_1$, da je $\det \phi(t, t_0) = 0$.
  Potem je matrika $\phi(t_1, t_0)$ singularna, zato ima netrivialno jedro,
  torej obstaja vsaj en neničeln vektor $d \in \R^n$, da $\phi(t_1, t_0) d = 0$.
  Torej
  \[
	\phi(t_1, t_0) d = \sum_{i=1}^n x_i(t_1) d_i = 0.
  \]
  Definirajmo
  \[
	z(t) = \sum_{i=1}^n d_i . x_i(t).
  \]
  Ta funkcija je rešitev sistema, zanjo velja $z(t_1) = 0$.
  Tudi funkcija $w(t) = 0$ je rešitev začetnega problema $\dot{x} = Ax, x(t_1) =
  0$, torej po eksistenčnem izreku $z = 0$.
  Ker so v točki $t_0$ vektorji $x_i$ linearno neodvisni, velja $d_i = 0$.
\end{proof}

\vprasanje{Kaj je fundamentalna matrika homogenega sistema linearnih NDE?
  Dokaži, da je nesingularna.}

Oglejmo si preslikavo $\mathcal{F}: (t_0-\varepsilon, t_0+\varepsilon) \times
\R^n \to \R^n$, definirano z
\[
  \mathcal{F}(t, y) = \phi(t, t_0) y.
\]
To je tok enačbe $\dot{x} = Ax$.
Za vsak dovolj majhen $t$ je preslikava $y \mapsto \mathcal{F}(t,y)$
difeomorfizem, saj je obrnljiva linearna preslikava $\R^n \to \R^n$.
Vzemimo $t_0 = 0$ in označimo $\phi(t,0) = \phi(t)$.

\begin{trditev}
  Velja $\phi(t_1 + t_2) = \phi(t_1) \phi(t_2)$.
\end{trditev}

\begin{proof}
  Po eni strani imamo za vsak $x \in \R^n$
  \begin{gather*}
	\phi(t_1) x = \mathcal{F}(t_1,x), \\
	\phi(t_2) \phi(t_1)x = \mathcal{F}(t_2, \phi(t_1) x) = \mathcal{F}(t_2 +
	t_1, x),
  \end{gather*}
  ker je $\mathcal{F}$ tok, po drugi strani pa
  \[
	\phi(t_1 + t_2) x = \mathcal{F}(t_1+t_2,x).
  \]
\end{proof}